\section{Loss of surface classes}\label{sec:cycles}

Once a target exceptional locus contains no surface, each surface in the
source exceptional locus forces a decrease in the rank of analytic
surface classes.  We prove this first for compact reduced analytic
spaces, so that it applies both to the fourfolds and to their possibly
nonnormal boundary components.  Restriction to the common open carries
source surface classes onto all target surface classes.  A removed
source surface lies in the kernel, and K\"ahler positivity makes its
class nonzero.
The use of surface-cycle ranks follows the algebraic argument of
Fujino \cite[Lemma~2.3]{Fujino2004}; we give the compact analytic form,
as in \cite[Section~4.1]{OrdinaryKahler}, with the additional generality
needed for these boundary components.

For a compact reduced complex analytic space $V$, define
\begin{equation}\label{eq:cycle-rank}
 \begin{split}
 \mathcal C_2(V)
 &:=\operatorname{span}_{\R}\{[S]\in H_4(V,\R):
       S\subset V\text{ is an irreducible compact analytic surface}\},\\
 c_2(V)&:=\dim_{\R}\mathcal C_2(V).
 \end{split}
\end{equation}
Here $[S]$ is the image in $H_4(V,\R)$ of the fundamental class of $S$,
with the complex orientation on its regular locus.  All homology groups
in this section have real coefficients.

We recall the topology that makes this definition and its use below
valid for singular analytic spaces.  A compact reduced analytic space
admits a finite triangulation compatible with any prescribed finite
collection of closed analytic subsets.  A subset of complex dimension
$d$ is triangulated in real dimension at most $2d$.  One can obtain this
statement for an abstract analytic space as follows.  Choose a compatible
Whitney stratification
\cite[Chapter~III, Propositions~1.5 and~2.2.2]{Teissier1982}, finitely
many analytic charts $z_j:U_j\hookrightarrow\mathbb C^{N_j}$, and nonnegative
smooth functions $\rho_j$ supported compactly in $U_j$, with positivity
sets covering the space.  Extending each block by zero gives an embedding
\[
 x\longmapsto\bigl(\rho_j(x),\rho_j(x)z_j(x)\bigr)_j
\]
into a Euclidean space.  Indeed, a block with $\rho_j(x)>0$ recovers its
chart coordinate as $w/t$.  Locally all blocks extend smoothly to an
ambient chart, and this left inverse shows that the embedding is locally
an ambient smooth embedding.  It therefore preserves the Whitney
conditions.  Mather's control data
\cite[Section~7, Proposition~7.1, and Section~8]{Mather1970} and
Goresky's triangulation theorem \cite[Section~5]{Goresky1978} give a
triangulation compatible with the strata and smooth on them.  Compactness
makes it finite; smoothness on the strata gives the stated dimension
bound.  In particular, $H_4(V,\R)$ is finite-dimensional and $c_2(V)$
is a finite nonnegative integer.

We will also use Borel--Moore homology $H_*^{\mathrm{BM}}$, defined by
locally finite chains.  It agrees with ordinary homology on compact
spaces, and it allows a fundamental class to be restricted to an open
subset even when that subset is noncompact.  For a closed analytic
subset $A\subset V$ of a compact analytic space, compatible
triangulations give the localization exact sequence
\begin{equation}\label{eq:bm-localization}
 \cdots\longrightarrow H_j(A,\R)\longrightarrow H_j(V,\R)
 \xrightarrow{\ j^*\ } H_j^{\mathrm{BM}}(V\setminus A,\R)
 \longrightarrow H_{j-1}(A,\R)\longrightarrow\cdots.
\end{equation}
This is the homological sequence for the closed subset and its open
complement.  In a compatible triangulation, the Borel--Moore group of
the complement is the homology of the relative chain complex for
$(V,A)$, so the sequence follows from the long exact sequence of that
pair; see
\cite[Sections~1.2 and~1.6]{BorelHaefliger1961}.
If $S$ is an irreducible surface, its singular
locus has real dimension at most two.  Applying this sequence to that
locus extends the oriented fundamental class of $S_{\mathrm{reg}}$
uniquely to a class in $H_4(S,\R)$.  This is the analytic fundamental
class used in \eqref{eq:cycle-rank}.  Its restriction to any open subset
is the corresponding Borel--Moore fundamental class; see also
\cite[Section~3.1 and Theorem~3.2]{BorelHaefliger1961}.

K\"ahler positivity supplies the nonvanishing needed for strict rank
loss.  If $V$ is a closed reduced analytic subspace of a compact
K\"ahler space $Y$, and $S\subset V$ is an irreducible surface, then
the local potentials of a K\"ahler form $\omega$ on $Y$ define a class
$[\omega]\in H^2(Y,\R)$.  Here is the construction on the possibly
singular space.  Let $\mathcal{PH}_Y$ be the sheaf of functions locally
equal to the real part of a holomorphic function.  The differences of
the potentials form a cocycle in this sheaf.  The connecting
homomorphism for the exact sequence
\[
 0\longrightarrow\R_Y\xrightarrow{\,\sqrt{-1}\,}\mathcal O_Y
   \xrightarrow{\,\operatorname{Re}\,}\mathcal{PH}_Y
   \longrightarrow0
\]
gives $[\omega]$, with the normalization agreeing with the de~Rham
class on smooth spaces.  Constant-sheaf cohomology agrees with
singular cohomology because analytic spaces are locally contractible,
as also follows from the compatible triangulations above.  The
construction commutes with holomorphic pullback; hence its pullback
to a resolution is the de~Rham class of the pulled-back form.  We obtain
\begin{equation}\label{eq:positive-surface-class}
 \bigl\langle[\omega]^2,[S]\bigr\rangle
   =\int_{S_{\mathrm{reg}}}\omega^2>0.
\end{equation}
The pairing here is taken in $Y$.  To compute it, choose a resolution
$r:\widetilde S\to S$ that is an isomorphism over $S_{\mathrm{reg}}$.
Since $r_*[\widetilde S]=[S]$, the pairing equals
$\int_{\widetilde S}(r^*\omega)^2$.  The pulled-back form is
semipositive everywhere and positive on the open set mapping to
$S_{\mathrm{reg}}$, which proves the strict inequality.  Consequently
the class of $S$ is nonzero in $H_4(Y,\R)$ and hence also in
$H_4(V,\R)$.

\begin{lemma}[Cycle loss]\label{lem:cycle-loss}
Let $P$ and $Q$ be compact reduced complex analytic spaces, and let
$L_P\subset P$ and $L_Q\subset Q$ be closed analytic subsets with
dense complements.  Suppose an isomorphism
\[
 U=P\setminus L_P\simeq Q\setminus L_Q
\]
is represented by a proper bimeromorphic graph: there is a compact
reduced analytic space $G$ with proper projections $p:G\to P$ and
$q:G\to Q$ such that the common open
\begin{equation}\label{eq:cycle-graph-open}
 p^{-1}(P\setminus L_P)=q^{-1}(Q\setminus L_Q)
\end{equation}
is dense in $G$ and maps isomorphically to both displayed complements.
If $\dim L_Q\leq1$, there is a surjection
\[
 \mathcal C_2(P)\longrightarrow\mathcal C_2(Q),
 \qquad\text{and thus}\qquad c_2(P)\geq c_2(Q).
\]
If $L_P$ contains an irreducible surface whose class in $H_4(P,\R)$ is nonzero,
then $c_2(P)>c_2(Q)$.  In particular, strict inequality holds whenever
$L_P$ contains a surface and $P$ is a closed analytic subspace of a
compact K\"ahler space.
\end{lemma}

\begin{proof}
The restriction from $Q$ will identify its fourth homology with the
fourth Borel--Moore homology of the common open.  We then compare the
surface classes by their restrictions there.

Since $L_Q$ has real dimension at most two,
$H_4(L_Q,\R)=H_3(L_Q,\R)=0$.  The degree-four part of
\eqref{eq:bm-localization} is therefore
\[
 0\longrightarrow H_4(Q,\R)
   \xrightarrow{\ j_Q^*\ }H_4^{\mathrm{BM}}(U,\R)
   \longrightarrow0.
\]
Define the linear map
\begin{equation}\label{eq:cycle-comparison-map}
 \Phi=(j_Q^*)^{-1}\circ j_P^*:
 H_4(P,\R)\longrightarrow H_4(Q,\R).
\end{equation}

To determine $\Phi$ on surface classes, we first check that strict
transforms are analytic, without a normality hypothesis.  Let
$S\subset P$ be an irreducible surface not contained in $L_P$.
Removing a proper analytic subset from an irreducible analytic space
leaves an irreducible space: its regular locus remains connected after
removing an analytic subset of positive complex codimension
\cite[Section~3.3]{BorelHaefliger1961}.  Thus
$S\cap U$ is irreducible.  The analytic inverse image $p^{-1}(S)$ has
a unique irreducible component $\widetilde S$ meeting the open where
$p$ and $q$ are isomorphisms; there it agrees with $S\cap U$.
This component is a compact irreducible surface.  Remmert's proper
mapping theorem shows that
\[
 S^+:=q(\widetilde S)\subset Q
\]
is an irreducible analytic surface, agreeing with $S\cap U$ on $U$.
It is the strict transform of $S$.  The same construction with $p$ and
$q$ interchanged gives strict transforms in the other direction.

The restrictions of $[S]$ and $[S^+]$ to $U$ are the same oriented
fundamental class, so \eqref{eq:cycle-comparison-map} gives
$\Phi([S])=[S^+]$.  If instead $S\subset L_P$, its restriction
vanishes and $\Phi([S])=0$.  It follows that $\Phi$ maps
$\mathcal C_2(P)$ into $\mathcal C_2(Q)$.  Every irreducible surface
in $Q$ meets $U$, since $\dim L_Q\leq1$; its strict transform in $P$
therefore maps back to its class.  The induced map of cycle spans is
surjective.  Finally, a nonzero surface class supported in $L_P$ lies
in its kernel.  Finite-dimensionality gives the strict rank inequality,
and \eqref{eq:positive-surface-class} gives the stated K\"ahler case.
\end{proof}

Here are the two applications to a flip
$T\to Z\leftarrow T^+$ with exceptional loci $L,L^+$.
Let $A\subset Z$ be the common exceptional image furnished by
Proposition~\ref{prop:comparison}.  On the normalized graph of the
flip, with projections $p,q$ and common composite $h$ to $Z$, one has
\[
 p^{-1}(T\setminus L)=h^{-1}(Z\setminus A)
                    =q^{-1}(T^+\setminus L^+).
\]
Both projections are isomorphisms there.  Thus this graph gives exactly
the diagram required by Lemma~\ref{lem:cycle-loss} for $P=T$, $Q=T^+$.
For a boundary component $H\subset T$ and its transform
$H^+\subset T^+$, give both spaces their reduced structures and put
\[
 L_H=H\cap L,\qquad L_{H^+}=H^+\cap L^+.
\]
Smallness ensures that the complementary opens are dense and
identified.  In the compact analytic fibre product $H\times_Z H^+$,
take the reduced irreducible component $G_H$ containing their common
open graph, with projections $p_H,q_H$ and common composite $h_H$ to
$Z$.  This component exists because the fibre product over
$Z\setminus A$ is precisely that graph.  Its projections are proper,
and the common open is dense in it.  Moreover
\[
 p_H^{-1}(H\setminus L_H)=h_H^{-1}(Z\setminus A)
                       =q_H^{-1}(H^+\setminus L_{H^+}),
\]
with both projections isomorphisms on this open.  This verifies
\eqref{eq:cycle-graph-open}, including equality of the inverse images;
normality of $H$ and $H^+$ is unnecessary.
In particular, when $H^+$ contains no flipped surface,
$\dim L_{H^+}\leq1$, so $c_2(H)\geq c_2(H^+)$.  The inequality is
strict if $H$ contains a flipping surface, because $H$ is a closed
analytic subspace of the K\"ahler fourfold $T$.  The surjectivity used
throughout is the surjectivity on the spans of surface classes proved
above; no surjectivity of \eqref{eq:cycle-comparison-map} on all homology
is needed.
